% =========================================================
% Proof by Contrapositive
% =========================================================
\paragraph{Proof by Contrapositive.}

The contrapositive of $P \Rightarrow Q$ is $\neg Q \Rightarrow \neg P$.
These two statements are logically equivalent: one is true if and only
if the other is. A proof by contrapositive establishes
$\neg Q \Rightarrow \neg P$ by a direct argument, and since this is
equivalent to $P \Rightarrow Q$, the original claim follows.

\begin{tcolorbox}[colback=propbox, colframe=propborder, arc=2pt,
  left=6pt, right=6pt, top=4pt, bottom=4pt,
  title={\small\textbf{Template: Proof by Contrapositive}},
  fonttitle=\small\bfseries]
\textbf{Goal.} $P \Rightarrow Q$.\\
\textbf{Equivalent goal.} $\neg Q \Rightarrow \neg P$.\\[4pt]
Assume $\neg Q$.\\
$\vdots$ \quad [derive $\neg P$ from $\neg Q$]\\
Therefore $\neg P$.\\
Hence $P \Rightarrow Q$. \qed
\end{tcolorbox}

\begin{remark}[When contrapositive is the right tool]
Contrapositive is useful when the negation of the conclusion is more
concrete and easier to work with than the hypothesis.

The canonical example is injectivity. To prove $P \Rightarrow Q$ where
$P$ says ``$f(a) = f(b)$'' and $Q$ says ``$a = b$'': the contrapositive
is $a \neq b \Rightarrow f(a) \neq f(b)$, which may be directly
accessible if the function is order-preserving or explicitly defined.

Another signal: when $Q$ is the statement that some object exists or
some set is non-empty. The negation $\neg Q$ says no such object
exists, which may be a very restrictive and workable condition.

The general diagnostic: after writing $\neg Q$, look at what it says.
If it gives you a strong, concrete assumption --- something you can
actually manipulate --- contrapositive is likely the right tool.
If $\neg Q$ is vague or complicated, try a different approach.
\end{remark}

\begin{remark}[Contrapositive versus contradiction: the essential difference]
Students routinely conflate contrapositive and contradiction.
They are not the same, and knowing the difference matters for clarity.

In a \textbf{contrapositive} proof, you assume \emph{only} $\neg Q$
and derive $\neg P$. The hypothesis $P$ is never mentioned.
The logical engine is: you proved $\neg Q \Rightarrow \neg P$,
which by equivalence gives $P \Rightarrow Q$.

In a \textbf{contradiction} proof, you assume \emph{both} $P$ and $\neg Q$
simultaneously and derive \emph{any} absurdity --- not necessarily $\neg P$.
The logical engine is: the assumption $P \land \neg Q$ leads to a false
statement, so $P \land \neg Q$ must itself be false, so $P \Rightarrow Q$.

When contrapositive works, it is cleaner: you need one assumption, not two,
and the conclusion you derive ($\neg P$) is specific rather than
``any contradiction.'' Always try contrapositive before reaching
for contradiction.
\end{remark}

\begin{example}[Contrapositive for divisibility]
\textbf{Claim.} If $n^2$ is even then $n$ is even.

\medskip
\noindent\textit{Proof (contrapositive).}
We prove the contrapositive: if $n$ is odd then $n^2$ is odd.

Assume $n$ is odd. Then $n = 2k+1$ for some $k \in \mathbb{Z}$.
\[
n^2 = (2k+1)^2 = 4k^2 + 4k + 1 = 2(2k^2+2k)+1,
\]
which has the form $2m + 1$ and is therefore odd.

Hence $n$ odd $\Rightarrow$ $n^2$ odd, which is the contrapositive of
$n^2$ even $\Rightarrow$ $n$ even. \qed

\medskip
\noindent
Notice: the hypothesis ``$n^2$ is even'' was never used. We proved
the contrapositive directly, and the original implication follows
by logical equivalence. There is no need to loop back and mention
the original statement; the contrapositive suffices.
\end{example}
